\section{Conclusion}
\label{chap:conclusion}

The report has described a methodology for structured pruning of a transformer encoder, applied to text classification tasks at high sparsity. The methodology combines a guided evolutionary architecture search under a hard compute budget, a surrogate distance metric scored on a small calibration set, and a multi phase distillation recovery procedure that restores task accuracy on the chosen architecture.

The conclusions the experimental evidence supports are modest. The methodology produces recovered accuracies competitive with the structured pruning literature at compute keep ratios between zero point three and zero point zero five on Question Natural Language Inference and Multi Genre Natural Language Inference. The evolutionary machinery contributes a measurable improvement over random and beam baselines. The cohort selection rule is the most consequential single design choice in the search, and the two phase structure is necessary for the recovery. The surrogate metric is a useful but imperfect predictor of recovered accuracy, with a Spearman rank correlation in absolute value around zero point five on the diagnostic set.

The methodology has known limitations that the report has documented. The surrogate's task term saturates at high sparsity. The trajectory term has a magnitude bias. The recovery procedure has known failure modes in its exponential moving average selection and in its train versus deploy gap at half skip layers. The architectures the search produces at the highest sparsity drop layers and starve the late layers in patterns that reflect both the structure of the surrogate metric and the structure of the task.

Future work that the present analysis points to includes a per layer reweighting of the trajectory term to compensate for its magnitude bias, a brief adaptation stage before the surrogate is computed to bring it closer to a measurement of recoverability, a modification of the deployment substitution rules to preserve layer normalisations at half skip layers, and an extension of the methodology to non classification tasks. Each of these directions is testable by a controlled experiment whose protocol the report's experimental sections provide a template for.

The methodology is, in our view, a reasonable point in the design space of structured pruning procedures. It is not a final answer. The components are individually well understood, the failure modes are documented, and the improvements remaining to be tested are concrete enough to be the subject of a follow up study.
